\section{One finite checker for a family of marker variants}\label{sec:markers}

\subsection{The family already present in the proposals}

M124 and FM215 are closely related members of the same predecessor-marker
family. An exact comparison finds that 429 of their 430 rows agree; their
only difference is
\[
 (127,1):\quad 1R59\quad\hbox{versus}\quad 1R109.
\]
Both providers satisfy the required zero-return equation:
\[
 (59,0)=(109,0)=(125,0)=0R62.
\]
The earlier M124 note explicitly permits providers 109 and 125. Accordingly,
we do not claim a new marker idea for the provider-125 instance. The new
contribution is a generic verified checker that packages the shared argument,
and its compiled application to three exact tables \cite{incoming}.

The crucial parent rows are
\begin{center}
\begin{tabular}{@{}crr@{}}\toprule
Old role & Read 0 & Read 1\\\midrule
124 & $0R62$ & $1R138$\\
127 & $0R14$ & $1R138$\\\bottomrule
\end{tabular}
\end{center}
Every old incoming row to role 124 writes a 1 and moves right; every old
incoming row to role 127 writes a 0 and moves right. Thus the bit immediately
left of the head distinguishes the two logical roles. The roles can be
identified at new control 124, while physical control 127 is reused to read
that origin marker.

\subsection{Generic rule and finite obligations}

Let the two old roles be \(a,b\), with a common one-successor \(w\) and
zero-successors \(u,v\). Define \(\phi(b)=a\) and \(\phi(q)=q\) otherwise.
The old zero transitions preserve zero and move right to \(u,v\); both old
one transitions preserve one and move right to \(w\). Let \(r_0,r_1\) be
existing return providers. The new local rows have the following shape:
\begin{align*}
 (a,0)&\mapsto 0Lb,&(a,1)&\mapsto 1R\phi(w),\\
 (b,0)&\mapsto 0Rr_0,&(b,1)&\mapsto 1Rr_1,\\
 (r_0,0)&\mapsto 0R\phi(v),&(r_1,0)&\mapsto 0R\phi(u).
\end{align*}
Outside \(a,b\), the new table agrees with the old table after mapping
source controls and target controls by \(\phi\).

The Boolean checker verifies these exact shapes, every ordinary row, both
incoming marker conditions, equal finite table lengths, distinct noninitial
roles, nonidentification with halt, and absence of outgoing halt rows. Its
soundness theorem covers indices beyond the stored tables as well: the
finite lookup returns \code{None} there, and equal lengths preserve that
behavior. No assumption about a particular startup or source register frame
is hidden in the concrete checks.

Here \(\Halt\) means reaching the designated halt label under the totalized
step function, which fixes a configuration on a missing row. The concrete
closure checks below identify this predicate with conventional halting.

\begin{theorem}[Generic checked marker fold]
If \code{check O N a b halt u v w r0 r1 = true}, then for every complete
Boolean tape \(t:\Z\to\{0,1\}\) and every head position \(h\),
\[
 \Halt(O,(0,h,t))\quad\Longleftrightarrow\quad\Halt(N,(0,h,t)).
\]
\end{theorem}
\begin{proof}
The old invariant consists of exactly two clauses:
\[
 q=a\Longrightarrow t(h-1)=1,
 \qquad
 q=b\Longrightarrow t(h-1)=0.
\]
It holds initially because \(a,b\ne0\). Each incoming-row condition
reestablishes its marker after one old step. An ordinary old transition is
matched by one new transition under \(\phi\). At either merged role on
read 1, the common one-successor gives the same one-step result. On read 0,
the new machine moves left to the selector, reads and preserves the origin
marker, moves right to the corresponding provider, then preserves the zero
and moves right to the required successor. The full tape is unchanged by
these three steps; the final head and mapped control match the old step.

Induction gives, for every old execution length \(n\), a new length \(m\)
with \(n\le m\le3n\) and identical tape/head at the matching boundary.
An old halt therefore transfers to a new halt. Conversely, suppose the new
machine halts by time \(n\). Simulate \(n\) old steps, reaching a matching
new time \(m\ge n\). Halt is absorbing in the totalized semantics, so the
new machine is still halted at \(m\). The checked equivalence
\(\phi(q)=\mathrm{halt}\iff q=\mathrm{halt}\) forces the old boundary
to be halted as well. This proves the reverse implication without another
local simulation or any fairness premise.
\end{proof}

The timing lemma is \code{timed\_simulation}; the checked finite endpoint
is \code{checked\_fold}. The generic source has 322 lines. The three
instances plus their explicit proposition-transfer corollary and audits
have 35 lines, with literal generated table data counted separately. This
is useful reuse, not a claim that the first generic implementation is
shorter than one preexisting hand proof.

\subsection{Concrete instances and M125 table}

All three instances set
\[
 (a,b,\mathrm{halt},u,v,w,r_0)=(124,127,254,62,14,138,8),
\]
and choose \(r_1\) as follows:
\begin{center}
\begin{tabular}{@{}llr@{}}\toprule
Instance & Provider $r_1$ & Working states\\\midrule
M124 & 59 & 215\\
FM215 & 109 & 215\\
M125, supplied here & 125 & 215\\\bottomrule
\end{tabular}
\end{center}

M125's exact four replacements relative to the primary native table are
\begin{center}
\begin{tabular}{@{}crr@{}}\toprule
Row & Primary K215 & M125\\\midrule
$110/0$ & $0R127$ & $0R124$\\
$124/0$ & $0R62$ & $0L127$\\
$127/0$ & $0R14$ & $0R8$\\
$127/1$ & $1R138$ & $1R125$\\\bottomrule
\end{tabular}
\end{center}
Its native SHA-256 is
\begin{center}\small\ttfamily
9359aa5ca6dee6db1ec26a32b47d9e95593c6de749ae53f6208ac4ef369256d3.
\end{center}
The archive includes both native and dense tables, the exact bijective
renaming, and their hashes. Each has 215 working controls and 430 rows,
with all working controls reachable in the transition graph. Graph
reachability is used for this structural census, not as a claim that all
controls occur in the blank computation.

The original native labels occupy 303 list slots in the formal data. The
unused slots are absent transitions, not working states. The finite audit
checks complete binary rows on the working controls and closure of every
target into the working controls or the designated halt. Starting from 0
on any tape therefore cannot reach a nonhalt missing-row slot. This relates
the totalized formal predicate, which means reaching label 254, to the
conventional halting interpretation of these exact machines.

\subsection{Fresh verification and precise scope}

The generic module, literal tables, and concrete instances compiled from
fresh sources. Six printed assumption reports are closed, and recursive
\code{coqchk -silent MarkerInstances} passed. A separate independent review
checked the finite checker and every literal row. The no-op-sharing tape
update is proved pointwise equal to ordinary tape update; no function
extensionality axiom is used.

An independent finite simulator additionally checks each synchronization
boundary through the complete calibrated blank prefix. For M124, FM215,
and M125 alike, 95,714 old steps correspond to 95,840 new steps, with exactly
63 three-step blocks. Both sides end at native control 84, head $-8$, and
the same complete tape containing 95 ones. These observations corroborate
the literal bindings. The arbitrary-tape theorem, not the finite prefix,
supplies the unbounded halting conclusion.

The compiled result is a native table-to-table theorem in this packet's
arbitrary-tape semantics. A formal bridge to the maintained zipper/dense
representation and composition with the full ZFC endpoint remain separate
integration tasks. Ordinary mathematical composition with the exact K215
calibration gives the same ZFC-formalizable consistency equivalence, but
that is not advertised as a freshly rebuilt complete assistant endpoint.

A corpus comparison found no rooted state-bijection or global-reflection
match for the explicit M125 table among 2,248 tracked \code{.tm} paths and
87 direct \code{.tm} members of the seven incoming archives. The qualifying
subset contains 640 complete numeric 215-state tables. One lacks an explicit
start header and is compared using start 0, as recorded. This is an inventory
fact about explicit tables. Because provider 125 was already described,
it supplies no discovery-priority claim for the construction.
